% ECONOMICS_PROBLEM: {"id": "G12-364", "title": "Climate risk versus green preferences", "jel_code": "G12", "source_id": "OP-0318"}
% !TeX program = xelatex
\documentclass[11pt,a4paper]{article}
\usepackage[margin=23mm]{geometry}
\usepackage{fontspec}
\setmainfont{Latin Modern Roman}
\usepackage{amsmath,amssymb,mathtools}
\usepackage{xcolor,enumitem,booktabs,longtable,array,xurl,hyperref}
\definecolor{muted}{HTML}{53636C}
\hypersetup{colorlinks=true,urlcolor=blue,linkcolor=blue}
\setlength{\parindent}{0pt}
\setlength{\parskip}{5pt}
\newcommand{\R}{\mathbb R}
\newcommand{\E}{\mathbb E}
\newcommand{\N}{\mathbb N}
\newcommand{\ind}{\mathbf 1}
\newcommand{\Var}{\operatorname{Var}}
\newcommand{\Cov}{\operatorname{Cov}}
\newcommand{\supp}{\operatorname{supp}}
\DeclareMathOperator*{\argmin}{arg\,min}
\DeclareMathOperator*{\argmax}{arg\,max}
\newcommand{\entrymeta}[3]{{\sffamily\small\color{muted}\textbf{\texttt{#1}}\quad|\quad #2\par #3\par}}
\newcommand{\Problemref}[1]{\texttt{#1}}
\newcommand{\JELref}[2]{\texttt{#2} (JEL \texttt{#1})}
\begin{document}
\section*{Climate risk versus green preferences}
\textbf{Problem ID:} \texttt{G12-364}\quad \textbf{JEL:} \texttt{G12}\par
% BEGIN ECONOMICS STATEMENT
\label{problem:OP-0318}
{\sffamily\small\color{muted}Original subject group: Asset pricing and financial markets\par}
\label{definitions:problem:30007600}
\textbf{Audit classification:} Published research agenda. The Bank agenda supports climate-risk pricing, while Adam and Nagel support subjective-risk research generally. Neither inspected passage establishes the exact green-taste/climate-risk decomposition. Price and return contributions must be kept in consistent units.
\subsection*{Definitions and precise research target}
Fix green and otherwise comparable conventional securities with payoff vector \(X_{t+1}\), positive price vector \(P_t\) and investor holdings. Gross returns are \(R^k_{t+1}=X^k_{t+1}/P^k_t\) for security \(k\). Investor \(i\) has ordinary consumption utility and a nonpecuniary preference parameter \(g_i\) for financing environmentally beneficial activity; investors may also hold different beliefs about climate states. Let \(Z_{t+1}\) be a specified physical or transition climate-risk factor and \(M_{i,t+1}\) the consumption discount factor. Within a separately specified taste-free common-belief benchmark, define climate risk compensation using covariance of returns with the benchmark pricing kernel; define a taste contribution as the equilibrium price change under \(do(g_i=0\text{ for all }i)\), keeping cash-flow technologies fixed. The task is to identify these distinct contributions to green valuation differences. Use stated assumptions linking surveys to beliefs and preferences, and variation that shifts tastes separately from climate exposures. Determine whether the joint law of prices, returns, holdings and survey responses distinguishes taste changes from anticipated climate cash-flow losses. If it does not, derive bounds or observational equivalence results. Account for firm responses when interpreting financing effects. A lower expected green return alone cannot identify either the climate-insurance motive or environmental preferences.

\textbf{Additional formal definitions and scope (editorial).} A taste model must specify how environmental benefit enters utility, for example \(u_i(c)+g_i G_i(q)\), where \(G_i\) is a stated benefit from holdings and \(g_i\geq0\); \(g_i\) is not the climate exposure of the security. Define physical and transition shocks and their payoff channels separately. In a taste-free, common-belief benchmark with an interior pricing kernel, \(E[R^k]-R^f=-\operatorname{Cov}(M,R^k)/E[M]\). With nonpecuniary benefits or constraints the ordinary consumption Euler equation has additional terms. The taste-price effect \(P(g)-P(0)\) is in price units, whereas a covariance risk premium is in return units: these are distinct estimands and are not an additive decomposition of one spread. To compare channels in common units, re-solve every counterfactual and use a declared price or return decomposition. Matching green and conventional claims requires their maturity, default exposure, taxes and liquidity to be stated.

For the identification part, declare an admissible class \(\Theta\) of structural models, including preferences, technologies, shock laws, measurement/sampling rules and any equilibrium selector. If \(O\) denotes the complete observable history and \(T(\theta)\) the target defined above, the identified set is \(\mathcal I_T=\{T(\theta):\theta\in\Theta,\ \mathcal L_\theta(O)=\mathcal L_{\rm obs}(O)\}\); point identification means this set is a singleton. A \(do(a)\) comparison replaces stated policy/technology equations while fixing the declared exogenous law and re-solving the remaining equations. These definitions do not assert that an identification theorem holds without further restrictions.
\subsection*{Variants and assumption changes}
\begin{enumerate}
\item \textbf{Climate risk only (editorial).} Set \(g_i=0\), impose common beliefs and an identified kernel, and estimate payoff climate exposures and risk compensation.
\item \textbf{Tastes and belief heterogeneity (editorial).} Permit \(g_i\) and subjective climate forecasts to vary independently; seek exclusion restrictions or bounds for taste-removal price effects.
\end{enumerate}
\subsection*{References and source audit}
\begin{enumerate}
\item Official January 9, 2026 web version; locator: The Financial System / System-wide efficiency, climate paragraph. Broad agenda; exact model editorial. URL: \url{https://www.bankofengland.co.uk/research/bank-of-england-agenda-for-research}.
\item Metadata and Section 6 verified (printed pp.35--36). Subjective-risk background; green-taste decomposition not stated there. URL: \url{https://bpb-us-w2.wpmucdn.com/voices.uchicago.edu/dist/f/575/files/2022/01/ExpectationsSurvey_13.pdf}.
\end{enumerate}
\begin{enumerate}\raggedright
\item\hypertarget{definitions:source:30007600:1}{} Bank of England. 2026. \emph{Bank of England Agenda for Research, 2025–2028}. The Financial System / System-wide efficiency, climate paragraph.
\url{https://www.bankofengland.co.uk/research/bank-of-england-agenda-for-research}.
\item\hypertarget{definitions:source:30007600:2}{} Klaus Adam; Stefan Nagel. 2022. \emph{Expectations Data in Asset Pricing}. Section 6, printed pp.35--36 (PDF pp.36--37).
\url{https://bpb-us-w2.wpmucdn.com/voices.uchicago.edu/dist/f/575/files/2022/01/ExpectationsSurvey_13.pdf}.
DOI: \url{https://doi.org/10.3386/w29977}.
\end{enumerate}

\subsection*{Common conventions: Source mathematical conventions}

These conventions apply to each entry unless explicitly replaced.

\textbf{Models and identification.} A primitive model \(m\) specifies all probability laws, feasible choices, information, equilibrium and measurement rules used by its target. The admissible class \(\mathcal M\) is fixed independently of outcomes. If \(P_O(m)\) is its observable probability law and \(\tau(m)\) the target, its identified set at observable law \(P\) is
\[
 \mathcal I_\tau(P)=\{\tau(m):m\in\mathcal M,\ P_O(m)=P\}.
\]
Point identification means that this set is a singleton. An empty set signals incompatibility of data and assumptions. Coordinate bounds need not describe the joint set. Bounds are sharp only if every retained target is attainable or arbitrarily approachable by compatible models. Finite bounds require explicit restrictions on unobserved quantities; unrestricted confounding, hidden wealth, extrapolation or arbitrary equilibrium selection can yield uninformative sets.

\textbf{Causal interventions.} A potential outcome \(Y_i(p)\) is the outcome under a complete policy regime \(p\), including timing, financing, information and market spillovers. The target population and units are fixed before treatment. With interference, use the complete assignment vector or a justified exposure mapping; individual no-interference notation alone cannot identify an economy-wide intervention. A difference of mean potential outcomes does not require the cross-world joint distribution. Conditioning on post-treatment migration, survival, enrollment or employment changes the estimand and needs separate justification.

\textbf{Statistical statements.} An estimator, confidence set, forecast or test specifies a sampling law, model class and loss/coverage criterion. Uniform \((1-\alpha)\)-coverage means \(\inf_{m\in\mathcal M}\Pr_m\{\tau(m)\in C_n\}\geq1-\alpha\), or an explicitly asymptotic version. The whole parameter space trivially has coverage, so informativeness requires a stated width, risk or power objective. A forecasting improvement is relative to a named benchmark, information set and evaluation population.

\textbf{Equilibrium and optimization.} An equilibrium comprises feasible plans, optimality, beliefs where needed, budget constraints and all market/resource clearing. The primitives or a stated benchmark define each correspondence \(\mathcal E(p;m)\). Nonemptiness is assumed or proved, never inferred just from compact policy sets. A selected-equilibrium target uses a declared selection rule; a robust target quantifies over all permitted equilibria. Use suprema or infima unless attainment is established. An \(\varepsilon\)-optimal policy is feasible and within \(\varepsilon\) of the finite optimum value. Report an empty feasible set separately.

\textbf{Welfare and accounting.} Normative utility, interpersonal normalization, social weights, numeraire, discounting and the use of public receipts are primitives. Behavioral utility need not coincide with the welfare criterion. Household welfare includes ownership income and rents once; adding profits again to compensating variation generally double-counts them. A monetary equivalent is defined by an explicit transfer and price convention, with monotonicity and range conditions. Average effects, mechanism contributions, transfers and welfare losses are different targets. Infinite sums require convergence and dynamic feasible plans require terminal or no-Ponzi conditions.

\textbf{Source versions.} A working-paper date, revision date and journal publication year are distinguished. A DOI for a journal version is not automatically the DOI of an earlier manuscript. Locators refer to the stated version and printed pages where specified. A metadata-only check does not verify a claim about a full-text conclusion. Every entry's source subsection narrows the provenance claim accordingly.
% END ECONOMICS STATEMENT
\end{document}
